\newpage
\appendix

\section{Evaluation Metrics}
\label{app:metrics}
To systematically assess model performance under controlled event-sequence generation, we design a human evaluation protocol based on five key aspects:

\begin{itemize}[leftmargin=*]
    \item \textbf{Narrative Quality (NQ)}: Evaluates how clearly and coherently the generated frames convey a complete story arc.
    \item \textbf{Subject Consistency (SC)}: Measures the consistency of key subject attributes across frames, including identity, appearance, and characteristics.
    \item \textbf{Prompt-Frame Alignment (PFA)}: Assesses how accurately the generated content reflects the input prompts in terms of scene, character, and emotional setting.
    \item \textbf{Shot Diversity (SD)}: Evaluates the richness and variability of camera angles and shot compositions throughout the sequence.
    \item \textbf{Scene Diversity (SceneD)}: Measures the variety and progression of environments and settings across frames, indicating spatial and temporal richness.
\end{itemize}

We employed a multi-round conversational evaluation approach, using three rounds of dialogue. In the first two rounds, we presented keyframes generated by two different models. The third round involved evaluating the results, with ratings assigned on a scale of 1-5. To mitigate any potential order effects, we evaluated the models' results in both forward and reverse order. If one model won both rounds, it was considered the winner; if each model won one round, the result was considered a tie. We evaluate our proposed metrics by systematically applying the structure outlined in Table \ref{tab:multi-turn prompt}, extracting and analyzing the corresponding results for each criterion.

% By using the system prompt: 

% "\textit{You are an expert in judgment and understanding based on frames. You will compare the frame descriptions generated by different models in historical dialogues according to specific requirements. Your analysis will focus on key elements in each dialogue round, such as the plot, characters, actions, and scenes, and you will assess the differences in the level of detail and coherence across these aspects.}".


\begin{longtable}[c]{c|>{\raggedright\arraybackslash}p{5.5cm}|>{\raggedright\arraybackslash}p{6cm}}
\toprule
\textbf{Round} & \textbf{User's Input} & \textbf{Assistant's Response} \\
\midrule
\endfirsthead

\midrule
\endfoot
1 & I provide the first set of story frames extracted from a long video. \newline \texttt{\textcolor{blue}{[Image 1]}} \texttt{\textcolor{blue}{[Image 2]}} ... & I accepted the first set of story frames generated by the first model. \\
\midrule
2 & I provide the second set of story frames from the same video. \newline \texttt{\textcolor{blue}{[Image 1]}} \texttt{\textcolor{blue}{[Image 2]}} ... & I accepted the second set of story frames generated by the second model. \\
\midrule
3 & \texttt{[Evaluation Criteria]} & \texttt{[Assistant's Response based on Evaluation Criteria]} \\
\label{tab:multi-turn prompt}
\end{longtable}


\section{Evaluation Criteria}

We present the evaluation criteria for all metrics, including Narrative Quality~(NQ), Subject Consistency~(SC), Prompt-Frame Alignment~(PFA), Shot Diversity~(ShotD), and Scene Diversity~(SceneD).


\begin{prompt}[title={Narrative Quality}, label=NQ]
0-1 point: There is no continuity between keyframes, making it impossible to discern the plot and lacking basic narrative logic.

1-2 points: There is a faint connection between keyframes, but the performance characteristics of each keyframe subject are the same, making it difficult to understand the overall story.

2-3 points: Keyframes can express the basic development of the story, but the scene transitions and story fluctuations still need improvement, and related elements are only described in a similar manner.

3-4 points: There is good continuity between keyframes, and the main story and plot development can be clearly depicted. There are rich descriptions and diverse corresponding presentation methods for the story development.

4-5 points: Keyframes form a perfect coherent narrative, accurately capturing all key moments and emotional changes, and having the rhythm of film narration.
\end{prompt}


\begin{prompt}[title={Subject Consistency}, label=SC]
0-1 point: The subject shows complete inconsistency across different keyframes, such as changes in appearance, clothing, and features, making it impossible to identify as the same subject.

1-2 points: The subject has obvious inconsistencies, but can still be recognized as the same subject, with significant changes in features affecting continuity.

2-3 points: The subject maintains basic consistency, with minor changes that do not affect overall recognition, and key features are generally consistent.

3-4 points: The subject maintains good consistency across all frames, with small and reasonable changes in features, clothing, and appearance, in line with the plot requirements.

4-5 points: The subject maintains perfect consistency across all keyframes, including uniform and natural details, expressions, and changes, meeting professional production standards.
\end{prompt}



\begin{prompt}[title={Prompt-Frame Alignment}, label=PFA]
0-1 point: The keyframes are almost irrelevant to the prompt, ignoring the core elements and scene descriptions in the prompt.

1-2 points: The keyframes have some relevance to the prompt, but only reflect a few descriptive elements, with most key information missing.

2-3 points: The keyframes basically conform to the prompt description, with the main elements reflected, but there are deviations in some details or emotional expressions from the prompt.

3-4 points: The keyframes are highly consistent with the prompt, accurately depicting the scene, characters, and emotions described, with slight differences in non-core details.

4-5 points: The keyframes perfectly interpret the content of the prompt, not only reflecting all explicit descriptions but also accurately capturing the implied atmosphere, style, and emotions.
\end{prompt}




\begin{prompt}[title={Shot Diversity}, label=ShotD]
0-1 point: All keyframes use a single, repetitive type of shot and angle, with monotonous composition and a lack of perspective variation.

1-2 points: Limited shot types, mainly using 2-3 basic shot angles, lacking in variation and creativity.

2-3 points: Some shot diversity, including several different types of shots (close-ups, medium shots, wide shots, etc.), but insufficient in innovation.

3-4 points: Rich shot types, diverse angles, and obvious composition changes, well-servicing narrative needs and emotional expression.

4-5 points: Demonstrates professional film-level shot language, including a variety of shot types, angles, and compositions, with strong innovation, with each shot precisely serving narrative and emotional expression.
\end{prompt}



\begin{prompt}[title={Scene Diversity}, label=SceneD]
0-1 point: All keyframes are confined to a single scene, environmental elements are extremely repetitive, and there is a lack of spatial and temporal changes.

1-2 points: Limited scene changes, mainly concentrated in 1-2 scenes, with minor changes in environmental elements and an unclear time span.

2-3 points: Shows several different scenes, with basic environmental changes, but scene transitions are not very natural or diversity is limited.

3-4 points: Scenarios are diverse and transitions are natural, including various environments (indoor/outdoor), with a clear time span and good use of space.

4-5 points: Scenarios are extremely rich and diverse, covering various types of environments, with excellent performance in time and space dimensions, and scene transitions perfectly match the narrative needs.
\end{prompt}

% \section{Broader Impact}
% Our work is developed using video data collected from publicly available sources and widely used open datasets. The dataset construction and model training pipelines include uniform sampling, deduplication, and structured prompt design to ensure content diversity and annotation neutrality. Annotations are generated using in-house or open large language models (e.g., Gemini), guided by structured instructions. No private or identifiable information is used. While our framework focuses on foundational research, appropriate usage guidelines and safeguards are expected to be followed to prevent unintended misuse. 
